\documentclass[10pt,a4paper]{article}

\usepackage[spanish,es-nodecimaldot]{babel}
\usepackage[utf8]{inputenc}
\usepackage[T1]{fontenc}
\usepackage{newtxtext,newtxmath}
\usepackage{geometry}
\usepackage{microtype}
\usepackage{amsmath,mathtools}
\usepackage{enumitem}
\usepackage{xcolor}
\usepackage{hyperref}
\usepackage{fancyhdr}
\usepackage[most]{tcolorbox}

\geometry{margin=2.5cm,includeheadfoot}
\setlength{\headheight}{15pt}
\setlength{\headsep}{0.28cm}
\setlength{\footskip}{0.62cm}
\setlength{\parindent}{0pt}
\setlength{\parskip}{0.28em}
\setlength{\abovedisplayskip}{0.5em}
\setlength{\belowdisplayskip}{0.5em}
\setlength{\abovedisplayshortskip}{0.25em}
\setlength{\belowdisplayshortskip}{0.25em}
\setlist[enumerate]{leftmargin=2.2em,itemsep=0.3em,topsep=0.25em}
\setlist[itemize]{leftmargin=1.5em,itemsep=0.2em,topsep=0.2em}

\definecolor{repblue}{gray}{0}
\definecolor{repteal}{gray}{0}
\definecolor{repgray}{gray}{0}
\definecolor{repaccent}{gray}{0}
\definecolor{repline}{gray}{0}
\definecolor{replight}{gray}{1}

\hypersetup{
  hidelinks,
  pdftitle={TRGF Padilla Tarea 1},
  pdfauthor={Juan Ignacio Padilla Barrientos}
}

\pagestyle{fancy}
\fancyhf{}
\fancyhead[L]{\textcolor{repblue}{\small\itshape TRGF Padilla}}
\fancyhead[R]{\textcolor{repgray}{\small Seminario UCR I 2026}}
\fancyfoot[C]{\textcolor{repgray}{\small \thepage}}
\renewcommand{\headrulewidth}{0.3pt}
\renewcommand{\footrulewidth}{0pt}

\newcommand{\CC}{\mathbb{C}}
\newcommand{\ZZ}{\mathbb{Z}}

\tcbset{
  enhanced,
  sharp corners,
  boxrule=0.45pt,
  arc=0pt,
  left=0.85em,
  right=0.85em,
  top=0.45em,
  bottom=0.4em,
  boxsep=0pt
}

\newcommand{\inlinehead}[2][repblue]{%
  \par\smallskip
  {\color{#1}\bfseries #2.}\hspace{0.45em}%
}

\newcommand{\blockhead}[2][repteal]{%
  \par\smallskip
  {\color{#1}\bfseries #2.}\par\nobreak\smallskip
}

\newcommand{\answer}{\inlinehead{Respuesta}}
\newcommand{\subhead}[1]{\blockhead{#1}}
\newcommand{\notehead}[1]{\inlinehead[repgray]{#1}}

\newcommand{\topbanner}{%
  \begin{tcolorbox}[
    colback=replight,
    colframe=repline,
    borderline west={1.2pt}{0pt}{repblue},
    borderline east={1.2pt}{0pt}{repteal},
    left=1.0em,
    right=1.0em,
    top=0.7em,
    bottom=0.55em
  ]
    {\Large\bfseries\color{repblue} TRGF Padilla \textemdash\ Tarea 1}\\[-0.1em]
    {\normalsize\color{repteal} Seminario UCR I 2026}
  \end{tcolorbox}
}

\newcounter{probcnt}
\newcommand{\problem}[2]{%
  \stepcounter{probcnt}%
  \ifnum\value{probcnt}>1\newpage\fi
  \par\vspace{0.25em}
  {\large\bfseries\color{repblue} Problema #1}\par
  {\color{repline}\rule{\linewidth}{0.55pt}}\par
  \begin{tcolorbox}[
    breakable,
    colback=white,
    colframe=repline,
    borderline west={1.6pt}{0pt}{repteal},
    left=0.7em,
    right=0.7em,
    top=0.35em,
    bottom=0.3em,
    before skip=0.35em,
    after skip=0.45em
  ]
    #2
  \end{tcolorbox}
}

\newcommand{\workspace}[1]{%
  \par\vspace*{#1\baselineskip}
}

\begin{document}

\color{black}
\topbanner
\small

\problem{1}{%
Calcule el determinante de grupo para $C_2 \times C_3$.
}

\answer
\begingroup
\fontsize{8.8}{10.2}\selectfont
\setlength{\abovedisplayskip}{0.28em}
\setlength{\belowdisplayskip}{0.28em}
\setlength{\abovedisplayshortskip}{0.18em}
\setlength{\belowdisplayshortskip}{0.18em}
\setlength{\arraycolsep}{3.2pt}
\renewcommand{\arraystretch}{0.96}

Ordenemos los elementos de
\[
  G=C_2 \times C_3=\langle a,b : a^2=e,\ b^3=e,\ ab=ba\rangle
\]
como
\[
  e,\ b,\ b^2,\ a,\ ab,\ ab^2.
\]
Escribimos $x_0:=x_e$, $x_1:=x_b$, $x_2:=x_{b^2}$, $y_0:=x_a$,
$y_1:=x_{ab}$, $y_2:=x_{ab^2}$. Como $a^{-1}=a$, $b^{-1}=b^2$,
$(b^2)^{-1}=b$, $(ab)^{-1}=ab^2$ y $(ab^2)^{-1}=ab$,
la matriz del determinante de grupo
\[
  \Theta_G=\det(x_{gh^{-1}})_{g,h\in G}
\]
es
\[
  M=
  \begin{pmatrix}
    x_0 & x_2 & x_1 & y_0 & y_2 & y_1\\
    x_1 & x_0 & x_2 & y_1 & y_0 & y_2\\
    x_2 & x_1 & x_0 & y_2 & y_1 & y_0\\
    y_0 & y_2 & y_1 & x_0 & x_2 & x_1\\
    y_1 & y_0 & y_2 & x_1 & x_0 & x_2\\
    y_2 & y_1 & y_0 & x_2 & x_1 & x_0
  \end{pmatrix}.
\]

\subhead{Reducción por bloques}
Escribamos
\[
  A=
  \begin{pmatrix}
    x_0 & x_2 & x_1\\
    x_1 & x_0 & x_2\\
    x_2 & x_1 & x_0
  \end{pmatrix},
  \qquad
  B=
  \begin{pmatrix}
    y_0 & y_2 & y_1\\
    y_1 & y_0 & y_2\\
    y_2 & y_1 & y_0
  \end{pmatrix},
\]
de modo que $M=\begin{pmatrix}A&B\\ B&A\end{pmatrix}$. Restando las tres
primeras filas de las tres últimas y luego sumando las tres últimas columnas a
las tres primeras, se obtiene
\begin{align*}
  \Theta_G
  &=\det(M)=\det(A+B)\det(A-B) \\
  &= \det
  \begin{pmatrix}
    x_0+y_0 & x_2+y_2 & x_1+y_1\\
    x_1+y_1 & x_0+y_0 & x_2+y_2\\
    x_2+y_2 & x_1+y_1 & x_0+y_0
  \end{pmatrix}
  \,\det
  \begin{pmatrix}
    x_0-y_0 & x_2-y_2 & x_1-y_1\\
    x_1-y_1 & x_0-y_0 & x_2-y_2\\
    x_2-y_2 & x_1-y_1 & x_0-y_0
  \end{pmatrix}.
\end{align*}

Para la matriz circulante
\[
  \det
  \begin{pmatrix}
    t_0 & t_2 & t_1\\
    t_1 & t_0 & t_2\\
    t_2 & t_1 & t_0
  \end{pmatrix}
  =(t_0+t_1+t_2)(t_0^2+t_1^2+t_2^2-t_0t_1-t_0t_2-t_1t_2).
\]
Aplicando esto a los dos determinantes anteriores, obtenemos
\begin{align*}
  \Theta_{C_2\times C_3}
  &= (x_0+x_1+x_2+y_0+y_1+y_2) \\
  &\quad\cdot \Bigl((x_0+y_0)^2+(x_1+y_1)^2+(x_2+y_2)^2 \\
  &\qquad\qquad -(x_0+y_0)(x_1+y_1)-(x_0+y_0)(x_2+y_2)-(x_1+y_1)(x_2+y_2)\Bigr) \\
  &\quad\cdot (x_0+x_1+x_2-y_0-y_1-y_2) \\
  &\quad\cdot \Bigl((x_0-y_0)^2+(x_1-y_1)^2+(x_2-y_2)^2 \\
  &\qquad\qquad -(x_0-y_0)(x_1-y_1)-(x_0-y_0)(x_2-y_2)-(x_1-y_1)(x_2-y_2)\Bigr).
\end{align*}
Esto calcula explícitamente el determinante de grupo de $C_2 \times C_3$.
\enlargethispage{2\baselineskip}
\endgroup

\problem{2}{%
Sea $G$ un grupo finito. Defina
\[
  [G,G] = \langle ghg^{-1}h^{-1} : g,h \in G \rangle,
\]
como el subgrupo de $G$ generado por todos los elementos de la forma
$ghg^{-1}h^{-1}$. Demuestre que $[G,G]$ es un subgrupo normal de $G$.
}

\answer
Escribimos $[g,h]:=ghg^{-1}h^{-1}$ para el conmutador de $g,h\in G$.
Debemos mostrar que $x[G,G]x^{-1}\subseteq [G,G]$ para todo $x\in G$.

Como $[G,G]$ está generado por conmutadores, todo elemento de $[G,G]$ es un
producto finito $w=[g_1,h_1]\cdots[g_n,h_n]$.
Observemos que
\[
  xwx^{-1}
  = (x[g_1,h_1]x^{-1})(x[g_2,h_2]x^{-1})\cdots(x[g_n,h_n]x^{-1}),
\]
pues insertar $x^{-1}x=e$ entre factores consecutivos no altera el producto.
Basta entonces verificar que el conjugado de cada conmutador es otro conmutador.

Sea $[g,h]=ghg^{-1}h^{-1}$ un conmutador arbitrario. Expandiendo:
\[
  x\,ghg^{-1}h^{-1}\,x^{-1}
  = (xgx^{-1})(xhx^{-1})(xg^{-1}x^{-1})(xh^{-1}x^{-1}).
\]
Usando que $xg^{-1}x^{-1}=(xgx^{-1})^{-1}$ y
$xh^{-1}x^{-1}=(xhx^{-1})^{-1}$, obtenemos
\[
  x[g,h]x^{-1}
  = (xgx^{-1})(xhx^{-1})(xgx^{-1})^{-1}(xhx^{-1})^{-1}
  = [xgx^{-1},\,xhx^{-1}].
\]
Como $xgx^{-1},xhx^{-1}\in G$, el lado derecho es un conmutador de elementos
de $G$ y por tanto pertenece a $[G,G]$.

Se sigue que $x[G,G]x^{-1}\subseteq [G,G]$ para todo $x\in G$, es decir,
$[G,G]\trianglelefteq G$.\hfill$\blacksquare$

\problem{3}{%
Demuestre que $G/[G,G]$ es un grupo conmutativo.
}

\answer
Escribimos $N:=[G,G]$ y $\bar g:=gN$ para la clase de $g$ en $G/N$.
Sean $g,h\in G$ arbitrarios. Debemos mostrar que $\bar g\bar h=\bar h\bar g$,
es decir, que $ghN=hgN$, lo cual equivale a $(hg)^{-1}(gh)\in N$.

Calculamos:
\[
  (hg)^{-1}(gh)=g^{-1}h^{-1}gh=[g^{-1},h^{-1}].
\]
Como $g^{-1},h^{-1}\in G$, este es un conmutador de elementos de $G$ y por
tanto pertenece a $[G,G]=N$.

Así, $\bar g\bar h=\bar h\bar g$ para todo $g,h\in G$, y $G/[G,G]$ es
conmutativo.\hfill$\blacksquare$

\problem{4}{%
Para un grupo $G$ no conmutativo, se define un carácter de la misma forma
que para el caso conmutativo; como un homomorfismo
\[
  \chi \colon G \to \CC^\times.
\]
Demuestre que si $\chi$ es un carácter de $G$, y $p$ es la proyección natural
\[
  p \colon G \to G/[G,G],
\]
entonces existe un carácter
\[
  \chi' \colon G/[G,G] \to \CC^\times
\]
tal que $\chi = \chi' \circ p$. (O sea, todo carácter se factoriza a través
de un carácter de un grupo conmutativo.)
}

\answer
Escribimos $N:=[G,G]$.

\subhead{$[G,G]\subseteq\ker\chi$}
Sea $[g,h]=ghg^{-1}h^{-1}$ un conmutador. Como $\chi$ es homomorfismo a
$\CC^\times$ y $\CC^\times$ es abeliano:
\[
  \chi([g,h])
  =\chi(g)\chi(h)\chi(g)^{-1}\chi(h)^{-1}
  =1.
\]
Así todo conmutador pertenece a $\ker\chi$, y como $\ker\chi$ es subgrupo,
$[G,G]\subseteq\ker\chi$.

\subhead{Definición de $\chi'$}
Definimos $\chi'\colon G/N\to\CC^\times$ por $\chi'(gN):=\chi(g)$.

\emph{Buena definición:} si $gN=hN$, entonces $h^{-1}g\in N\subseteq\ker\chi$,
luego $\chi(h^{-1}g)=1$, es decir, $\chi(g)=\chi(h)$.

\emph{Homomorfismo:} $\chi'(gN\cdot hN)=\chi'(ghN)=\chi(gh)=\chi(g)\chi(h)
=\chi'(gN)\chi'(hN)$.

Finalmente, $(\chi'\circ p)(g)=\chi'(gN)=\chi(g)$ para todo $g\in G$,
así que $\chi=\chi'\circ p$.\hfill$\blacksquare$

\problem{5}{%
Encuentre todos los caracteres de $S_3$.
}

\answer
Por el Problema~4, todo carácter $\chi\colon S_3\to\CC^\times$ se
factoriza a través de $S_3/[S_3,S_3]$, así que basta encontrar
$[S_3,S_3]$ y luego los caracteres del cociente.

\subhead{Cálculo de $[S_3,S_3]$}
Computamos un conmutador:
\[
  [(12),(13)]
  = (12)(13)(12)(13)
  = (132)(12)(13)
  = (23)(13)
  = (123).
\]
Como $(132)=(123)^{-1}$ y $\langle(123)\rangle = A_3 = \{e,(123),(132)\}$,
se tiene $A_3\subseteq [S_3,S_3]$.

Para la otra contención, $S_3/A_3\cong C_2$ es abeliano.
Si $N\trianglelefteq G$ y $G/N$ es abeliano, entonces para todo $g,h\in G$
se tiene $ghN=hgN$, es decir $(hg)^{-1}gh\in N$, o sea $[g,h]\in N$;
luego $[G,G]\subseteq N$.
Aplicando esto con $N=A_3$, obtenemos $[S_3,S_3]\subseteq A_3$.

Por lo tanto $[S_3,S_3]=A_3$.

\subhead{Caracteres}
El cociente $S_3/A_3\cong C_2=\{1,-1\}$ tiene exactamente dos caracteres, que
al componerse con la proyección $p\colon S_3\to S_3/A_3$ dan:
\begin{itemize}
  \item $\chi_0\colon S_3\to\CC^\times$,\; $\chi_0(\sigma)=1$ para toda
        $\sigma$ \textup{(carácter trivial)}.
  \item $\chi_1\colon S_3\to\CC^\times$,\; $\chi_1(\sigma)=\operatorname{sgn}(\sigma)$
        \textup{(carácter signo):} vale $+1$ en $A_3$ y $-1$ en $S_3\setminus A_3$.
\end{itemize}
Estos son todos los caracteres de $S_3$.\hfill$\blacksquare$

\problem{6}{%
(*) Encuentre y factorice el determinante de grupo de $S_3$.
}

\answer
\begingroup
\fontsize{8.8}{10.2}\selectfont
\setlength{\abovedisplayskip}{0.28em}
\setlength{\belowdisplayskip}{0.28em}
\setlength{\abovedisplayshortskip}{0.18em}
\setlength{\belowdisplayshortskip}{0.18em}
\setlength{\arraycolsep}{3.2pt}
\renewcommand{\arraystretch}{0.96}

Ordenamos $S_3$ como $e,(123),(132),(12),(13),(23)$ (primero~$A_3$, luego
su clase lateral) y abreviamos
\[
  a:=x_e,\; b:=x_{(123)},\; c:=x_{(132)},\;
  d:=x_{(12)},\; p:=x_{(13)},\; q:=x_{(23)}.
\]

\subhead{Estructura de bloques}
Calculando $g_ih_j^{-1}$ para cada par, la matriz
$M=(x_{g_ih_j^{-1}})$ resulta $M=\bigl(\begin{smallmatrix}A&B\\B&A\end{smallmatrix}\bigr)$ con
\[
  A=\begin{pmatrix}a&c&b\\b&a&c\\c&b&a\end{pmatrix},\qquad
  B=\begin{pmatrix}d&p&q\\p&q&d\\q&d&p\end{pmatrix}.
\]
La igualdad de los bloques fuera de la diagonal se verifica usando que
conjugar $A_3$ por $(12)$ intercambia $(123)\leftrightarrow(132)$;
la del bloque inferior derecho, porque el producto de dos transposiciones
distintas es una rotación, así que ese bloque reproduce la tabla de~$A_3$.
Como en el Problema~1:
\[
  \Theta_{S_3}=\det M=\det(A+B)\,\det(A-B).
\]

\subhead{Factores lineales}
Cada fila de $A+B$ suma $L_+:=a{+}b{+}c{+}d{+}p{+}q$
y cada fila de $A-B$ suma $L_-:=a{+}b{+}c{-}d{-}p{-}q$,
así que $(1,1,1)^{T}$ es autovector de ambas matrices.
Sumando las tres filas y restando la primera columna de las demás:
\[
  \det(A+B)=L_+\cdot\det
  \begin{pmatrix}
    a{-}b{-}p{+}q & c{-}b{+}d{-}p\\
    b{-}c{+}d{-}q & a{-}c{+}p{-}q
  \end{pmatrix},
\]
y análogamente $\det(A-B)=L_-$ multiplicado por el determinante~$2\times 2$
que se obtiene del anterior al sustituir $(d,p,q)\to-(d,p,q)$.

\subhead{Factor cuadrático}
Denotemos las entradas del determinante~$2\times 2$ de $A+B$ como
\[
  \det\begin{pmatrix}\alpha&\beta\\\gamma&\delta\end{pmatrix}
  =\alpha\delta-\beta\gamma,
\]
con $\alpha=a{-}b{-}p{+}q$, $\beta=c{-}b{+}d{-}p$,
$\gamma=b{-}c{+}d{-}q$, $\delta=a{-}c{+}p{-}q$.
Expandimos cada producto por separado.

Producto diagonal:
\begin{align*}
  \alpha\delta
  &= (a{-}b{-}p{+}q)(a{-}c{+}p{-}q)\\
  &= a^2-ac+ap-aq-ab+bc-bp+bq\\
  &\quad -ap+cp-p^2+pq+aq-cq+pq-q^2\\
  &= a^2-ab-ac+bc-p^2-q^2+2pq+cp-bp-cq+bq.
\end{align*}
Producto antidiagonal:
\begin{align*}
  \beta\gamma
  &= (c{-}b{+}d{-}p)(b{-}c{+}d{-}q)\\
  &= cb-c^2+cd-cq-b^2+bc-bd+bq\\
  &\quad +db-dc+d^2-dq-pb+pc-pd+pq\\
  &= -b^2-c^2+2bc+d^2-cq+bq-dq-pb+pc-pd+pq.
\end{align*}

Restando:
\begin{align*}
  Q = \alpha\delta-\beta\gamma
  &= (a^2+b^2+c^2-ab-ac-bc)\\
  &\quad -(d^2+p^2+q^2-dp-dq-pq).
\end{align*}
Cada monomio del segundo grupo ($d^2,p^2,q^2,dp,dq,pq$) tiene grado par
en las variables $d,p,q$, es decir, es invariante bajo el cambio
$(d,p,q)\mapsto(-d,-p,-q)$. En consecuencia, al sustituir
$(d,p,q)\to-(d,p,q)$ en la matriz~$2\times 2$ de $A-B$, su determinante
vale también~$Q$.

\subhead{Resultado}
\[
  \boxed{\;\Theta_{S_3}
  = L_+\;\cdot\; L_-\;\cdot\; Q^{\,2}\;}
\]
con $L_\pm=a{+}b{+}c\pm(d{+}p{+}q)$ y
$Q=(a^2{+}b^2{+}c^2{-}ab{-}ac{-}bc)-(d^2{+}p^2{+}q^2{-}dp{-}dq{-}pq)$.
\hfill$\blacksquare$
\enlargethispage{2\baselineskip}
\endgroup

\end{document}
